\documentclass[10pt,a4paper]{amsart}
\usepackage[margin=19mm]{geometry}
\usepackage{amsmath,amssymb,mathtools}
\usepackage[scr=rsfs,cal=euler]{mathalfa}
\usepackage{xfrac}
\usepackage[unicode,hidelinks,bookmarks=false]{hyperref}
\newcommand{\ie}{i.e.\ }
\newcommand{\cf}{cf.\ }
\newcommand{\resp}{respectively\ }
\newcommand{\Subsection}{\subsection}
\providecommand{\nobreakdash}{\nobreak\hbox{-}}
\setlength{\emergencystretch}{2em}
\multlinegap0pt
\usepackage{amssymb}
\usepackage{tikz}
\usepackage{tensor}
\usepackage{xcolor}
\newtheorem{thm}{Theorem}
\newcommand{\ii}{\textnormal{i}}
\title{\bf{Non-perturbative topological strings\\ from resurgence}}
\author{Murad Alim}
\date{Source: arXiv:2406.17852v2 (CC BY 4.0)}
\begin{document}
\maketitle

\noindent\emph{Source and licence.} \bf{Non-perturbative topological strings\\ from resurgence}. Version arXiv:2406.17852v2 (CC BY 4.0), distributed by arXiv under the Creative Commons Attribution 4.0 International licence, http://creativecommons.org/licenses/by/4.0/, as recorded in the arXiv licence metadata for this exact version and shown on the abstract page https://arxiv.org/abs/2406.17852v2. Full licence notice: \texttt{licenses/arxiv-2406.17852-CC-BY-4.0.txt}.\par\smallskip

\noindent\textit{Editorial scope.} Counted unit: the article's thm thm:Borel (source0.tex lines 1047--1064). The article's complete original proof of it is delivered with it (source0.tex lines 1066--1115, one \texttt{proof} environment of 50 lines). No mathematics is written, repaired or re-derived: the statement and the proof are the article's own lines, in the article's own notation, and no retained line is substituted or re-flowed.  Retained uncounted as the source-local context the proof needs: lines 984--987.  Nothing was omitted from any retained span, so the package carries no omission record.  No retained line contains a citation, so the wrapper carries no bibliography and no bibliography entry.  No retained line contains a figure environment, an image-inclusion command or drawing code, so no image is delivered and the package declares no asset dependency.  Editorial formatting only, in addition: an AMS \texttt{amsart} preamble that restates the article's own declarations and macros used by the retained lines, namely the environments \texttt{thm} on separate counters, together with the standard packages those lines need.  The retained environments print in this document as Theorem 1 in order of appearance, whereas the article prints the same items with the numbering of its own full document.  The complete native source of the article as distributed in the arXiv e-print for this exact version is bundled unchanged as \texttt{sources/161378/source0.tex}.
\begin{align}\label{eq:symBern}
&\frac{e^{w(x+1)}+e^{w(-x+1)}}{ (e^w-1)^2} = \frac{2}{w^2}+ \left( -\frac{1}{6}+x^2\right)  \nonumber \\ 
&- \sum_{g=2}^{\infty}  \frac{1}{2g!} \left(  (2g-1) (B_{2g}(x)+B_{2g}(-x))- 2g x (B_{2g-1}(x)-B_{2g-1}(-x))\right)  w^{2g-2}\,.
\end{align}

\begin{thm} \label{thm:Borel}
Let $\Phi(\check{\lambda},t):= \sum_{g=2}^{\infty} \lambda^{2g-2} \tilde{F}^g(t)=  \sum_{g=2}^{\infty} \check{\lambda}^{2g-2} \, (2\pi)^{2g-2} \tilde{F}^g(t)\,, \check{\lambda}=\lambda/2\pi$ denote the asymptotic series of the generating function of  higher genus $g \ge 2$ GW invariants. Let furthermore the formal variables $t^{\beta} \in \mathbb{C}^{\times}\setminus \mathbb{Z}$ be with $\textrm{Re}(t^\beta) <1/2$ for all $\beta \in H_2(X,\mathbb{Z})$, then the Borel transform $G(\xi,t)$ of $\Phi(\check{\lambda},t)$  can be analytically continued to a meromorphic function given by:
\begin{equation} \label{eq:BorelTransform}
G(\xi,t) 
=  \frac{1}{\xi} \sum_{\beta>0} \sum_{k,n\in \mathbb{Z}} \frac{\Omega_{\beta, n}}{2}
\left(
  \frac{2(t^{\beta}+k)^2}{\xi^2}
  + \left(n^2 - \frac{1}{6}\right)
  - \frac{
      e^{\frac{\xi(n+1)}{t^{\beta}+k}}
      + e^{\frac{\xi(-n+1)}{t^{\beta}+k}}
    }{
       \left( e^{\frac{\xi}{t^{\beta}+k}} - 1 \right)^2
    } 
\right)\,,
\end{equation}
with poles at $\xi= 2\pi \ii m  (t^\beta+k)\,,$ for $m\in \mathbb{Z}\setminus \{0\}$ and $k\in \mathbb{Z}$.
\end{thm}

\begin{proof}

The Borel transform is defined as the formal power series $G(\xi,t):=\mathcal{B}(\Phi(-,t))(\xi)$, where
\begin{equation}
    \mathcal{B}\colon\check{\lambda}\mathbb{C}[[\check{\lambda}]]\to \mathbb{C}[[\xi]], \;\;\;\; \mathcal{B}(\check{\lambda}^{n+1})=\frac{\xi^n}{n!}.
\end{equation}
We wish to study the Borel transform of 
\begin{align}
\Phi(\check{\lambda},t)& = \sum_{g=2}^{\infty}\sum_{\beta>0} \sum_{n\in\mathbb{Z}}\frac{(2\pi \ii)^{2g-2} \Omega_{\beta,n}}{2\cdot2g!} \left( (2g-1) (B_{2g}(n)+B_{2g}(-n))- 2g n (B_{2g-1}(n)-B_{2g-1}(-n))\right)\nonumber \\
&\times  \operatorname{Li}_{3-2g}(Q^{\beta}) \check{\lambda}^{2g-2}\,,
\end{align}
which is given by
\begin{align}
G(\xi,t)& = \sum_{g=2}^{\infty}\sum_{\beta>0} \sum_{n\in\mathbb{Z}}\frac{(2\pi \ii)^{2g-2} \Omega_{\beta,n}}{2\cdot2g!} \left( (2g-1) (B_{2g}(n)+B_{2g}(-n))- 2g n (B_{2g-1}(n)-B_{2g-1}(-n))\right)\nonumber \\
&\times  \operatorname{Li}_{3-2g}(Q^{\beta}) \frac{\xi^{2g-3}}{2g-3!}\,,
\end{align}

we now use the identity
$$
\operatorname{Li}_s \left(e^{\mu}\right)
=
\Gamma(1 - s)
\sum_{k=-\infty}^{\infty}
\left(2k\pi i - \mu\right)^{\,s - 1}
$$

$$
\text{valid for } \textrm{Re}(s) < 0 \quad \text{and all } \mu \text{ such that } e^{\mu} \neq 1,
$$
to write

\begin{align}
G(\xi,t)& = \frac{1}{\xi} \sum_{g=2}^{\infty}\sum_{\beta>0} \sum_{k\in \mathbb{Z}}\sum_{n\in\mathbb{Z}} \frac{\Omega_{\beta,n}}{2\cdot (2g)!} \left( (2g-1) (B_{2g}(n)+B_{2g}(-n))- 2g n (B_{2g-1}(n)-B_{2g-1}(-n))\right)\nonumber \\
& \times \left(\frac{\xi}{t^{\beta}+k}\right)^{2g-2}\,,
\end{align}
we can now use \eqref{eq:symBern} to write this as:
\begin{equation}
G(\xi,t)= \frac{1}{\xi} \sum_{\beta>0} \sum_{k,n\in \mathbb{Z}} \frac{\Omega_{\beta, n}}{2}
\left(
  \frac{2(t^{\beta}+k)^2}{\xi^2}
  + \left(n^2 - \frac{1}{6}\right)
  - \frac{
      e^{\frac{\xi(n+1)}{t^{\beta}+k}}
      + e^{\frac{\xi(-n+1)}{t^{\beta}+k}}
    }{
       \left( e^{\frac{\xi}{t^{\beta}+k}} - 1 \right)^2
    } 
\right)\,.
\end{equation}
\end{proof}

\end{document}
